\clearpage\SourceLeaf{517}
\renewcommand{\BookRunningTitle}{Hülfstafel bei Auflösung der unbestimmten Gleichung}
\pdfbookmark[0]{Hülfstafel bei Auflösung der unbestimmten Gleichung}{hulfstafel-ausschliessung}
\begin{center}
\textbf{HÜLFSTAFEL}\\[.7em]
BEI AUFLÖSUNG DER UNBESTIMMTEN GLEICHUNG
\[
A=fxx+gyy
\]
VERMITTELST DER AUSSCHLIESSUNGSMETHODE.
\end{center}
\BookRule{.18\textwidth}
Es wird vorausgesetzt, dass man zum Excludens eine Primzahl $p$ gewählt habe, durch welche keine der Zahlen $A$, $f$, $g$ theilbar ist. Auch beschränkt sich die Tafel auf die zwei Fälle, da der Werth des Ausdrucks $\frac{A}{f}(\mathrm{mod}.p)$ entweder ein bestimmter quadratischer Rest (allemal $1$), oder ein bestimmter quadratischer Nichtrest des Modulus $p$ ist. Endlich hat man sich begnügt, die Tafel nur für den Fall einzurichten, wo $fg$ ein quadratischer Rest von $p$ ist, und den entgegengesetzten ganz übergangen. Der sechste Abschnitt der \textit{Disquisitiones Arithmeticae} gibt hinlängliche Belehrung, wie man das, was die Tafel nicht unmittelbar enthält, leicht aus derselben ableiten könne.
\par\medskip\textit{Beispiele.} Es sei die aufzulösende Gleichung $21680143=xx+78yy$
\par\medskip\noindent 1) Excludens $=5$\hfill $fg=\mathrm{N},\quad \frac{A}{f}\equiv3\equiv2\pdot2^2$
\par\noindent Ex tabula $1,4$ pro $fg=R$ adeoque $0,2,3$ pro $fg=\mathrm{N}$,
et pro cas upr. $0,1,4$ sive excl. $5n\pm2$
\par\medskip\noindent 2) Excludens $=7$\hfill $fg=R,\quad \frac{A}{f}\equiv2\equiv1,3^2$
\par\noindent Ex tabula $0,1,2,5,6$ Pro casu praes. $0,3,6,1,4$ et excl. $7n\pm2$
\par\medskip\noindent 3) Excludens $=11$\hfill $fg=R,\quad \frac{A}{f}\equiv1$
\par\noindent habentur itaque $0,1,3,5,6,8,10$ excl. $11n\pm2,4$
\par\medskip\noindent 4) Excludens $=17$\hfill $fg=\mathrm{N},\quad \frac{A}{f}\equiv9\equiv1\pdot3^1$
\[
\begin{gathered}
0\pdot1\pdot3\pdot4\pdot6\\
0\pdot3\pdot8\pdot5\pdot1\\
0,2,3,4,6,7\\
\text{Excludens}\quad\pm1,5,8
\end{gathered}
\]
